\documentclass[11pt,a4paper]{article}
\usepackage[utf8]{inputenc}
\usepackage[T1]{fontenc}
\usepackage[margin=2.4cm]{geometry}
\usepackage{amsmath,amssymb,booktabs,tabularx,array,enumitem,hyperref,xcolor,listings}
\newcommand{\red}[1]{\textcolor{red}{#1}}
\newcommand{\blue}[1]{\textcolor{blue}{#1}}
\newcommand{\green}[1]{\textcolor{green}{#1}}
% Mavi paragraflar: ilk aşamada yaptıklarımız ve ana bulgular (Kasım 2026).
\newenvironment{progress}{\par\smallskip\color{blue}}{\par\smallskip}

\hypersetup{colorlinks=true,linkcolor=blue,urlcolor=blue}
\lstdefinestyle{cpp}{language=C++,basicstyle=\ttfamily\footnotesize,breaklines=true,frame=single,columns=fullflexible,showstringspaces=false}
\title{Measurement-Constrained Generation of a Geographical LTE Network in ns-3\\\large Technical Methodology and Implementation Plan (Working Draft)}
\author{Research working document}
\date{November 2026 --- updated S\i nd\i rg\i{} case study}
\begin{document}\maketitle
\begin{abstract}
This document specifies a reproducible procedure to construct a \emph{representative}, geographically anchored LTE network in ns-3 using Turkcell cell identifiers, shared site coordinates, and non-contemporaneous cell-level radio and traffic measurements. The objective is not to replay the operator network exactly. Unknown radio configurations and user locations are represented by explicit assumptions and scenario ensembles. The resulting simulator will support later studies of coverage, load, service deficits, controlled outages, and aerial network restoration.
\end{abstract}

\begin{progress}
Mavi notlar, ilk aşamada yaptıklarımızı ve ana bulguları ilgili bölümlerin altında özetliyor. Kısaca: Sındırgı ağı ns-3'te çalışıyor ve gerçek ölçümlerle iyi örtüşüyor. Ayarları seçerken kullanmadığımız üç kontrol sahasında, simülasyon ile Experia arasındaki ortalama fark sinyal gücü (RSRP) için 1,7~dB, sinyal kalitesi (SINR) için 2,2~dB. Her sonucu üç farklı anten yönü ve kullanıcı yerleşimiyle tekrarlayarak kontrol ettik. Kodlar ve ayrıntılı raporlar proje klasöründe duruyor.
\end{progress}
\section{Scope and key methodological constraints}
The available data include (i) cell identifiers and coordinates (several cells share each physical site coordinate), (ii) aggregate radio-quality indicators such as RSRP, RSRQ, SINR, and mean timing advance (TAAVG), and (iii) a separate LTE KPI dataset with traffic, users, PRB utilization, CQI, throughput, and potentially availability, drops, and handover counters. The radio and traffic observation periods \textbf{do not coincide}, although cell IDs and site coordinates match. Thus, these datasets may be joined by cell identity for metadata and long-run characterization, \emph{not} treated as simultaneous observations. Actual carrier frequencies, bandwidths, antenna azimuths, transmit powers, schedulers, and user coordinates are unavailable.

The deliverable is a \emph{measurement-constrained scenario generator}, not a digital replica of the operator deployment. Claims about actual coverage boundaries, actual failure locations, and measured traffic demand must be avoided without independent evidence.

\section{Inputs and data audit}
Prepare two normalized tables and a site index:
\begin{itemize}[nosep]
\item \texttt{cells.csv}: cell\_id,site\_id,lat,lon,radio\_median\_rsrp,radio\_median\_rsrq,radio\_median\_sinr,taavg,...; retain measurement period and units.
\item \texttt{kpis.csv}: timestamp,cell\_id,dl\_bytes,ul\_bytes,active\_users,dl\_prb,ul\_prb,cqi,throughput,...; retain counter versus rate semantics and time interval.
\item \texttt{sites.csv}: one row per unique site coordinate, listing its constituent cell IDs.
\end{itemize}
\textbf{Checklist:} (1) verify one-to-one cell-ID linkage across sources; (2) count cells per unique coordinate; (3) check coordinate precision and accidental near-duplicates; (4) check whether cell IDs encode sector or carrier information; (5) inspect missingness and KPI units; (6) determine whether TAAVG is in LTE TA commands, meters, or vendor-specific units; (7) identify whether throughput is aggregate or per-user; (8) distinguish carried traffic from offered demand. If the number of cells per site exceeds three, do not automatically create that many antenna directions: some cells may be co-sector carriers.

\begin{progress}
Experia ve KPI verileri hücre kimliği üzerinden birbirine bağlanabiliyor. KPI'daki saha adı sütunu her satırda aynı olduğu için bu iş için kullanılamıyor. İki veri farklı dönemlere ait (Experia 2026, KPI 2025); bu yüzden onları aynı anın ölçümü gibi birleştirmedik. TAAVG, throughput ve trafik hacminin birimleri belli olmadığı için bunları hesaplarda kullanmadık. Ağın yükünü bunun yerine birimsiz olan kapasite kullanım oranıyla (PRB) tanımladık.
\end{progress}

\section{Geographical topology generation}
\red{ben asagiya bir region sectim. deprem merkezine yakin, AIya koordinatlari girdim ona gore sitelari buldu ama kontrol edin}

\begin{progress}
Seçtiğiniz 16 sahanın konumlarını kontrol ettik, hepsi doğru. İlk koordinat dosyasında eksik olan 553, 554 ve 555 güncellenen dosyada var ve konumları seçiminizle aynı. Sınır etkisini görmek için ana 8 sahayı sabit tutup komşu saha sayısını artırdık (toplam 16, 24, 32 ve 48 saha). 24 sahadan sonra ana sahalardaki sonuçlar neredeyse hiç değişmiyor. Bu yüzden ağı sizin 16 sahanız ve en yakın 8 sahayla, toplam 24 sahada sabitledik.
\end{progress}
\subsection{Choose a bounded test region}
Select a contiguous cluster of approximately 10--20 physical sites, with surrounding sites included as an interference/association buffer where possible. Report the region boundary, number of physical sites, number of logical cells, inter-site distance distribution, and the fraction of sites near the simulation boundary. Avoid interpreting edge artifacts as network vulnerabilities.

\subsection{Convert geographical coordinates}
Use a local projected coordinate system (e.g., UTM appropriate to the region, or a local east--north tangent-plane approximation for a compact area) to obtain coordinates $(x_i,y_i)$ in meters. Translate the origin for numerical convenience. Store the coordinate transformation and original site IDs so maps can be reproduced.

\subsection{Distinguish sites, sectors, and carriers}
\red{bir cok siteda birden fazla cell var hepsinin konumu ayni gorunuyor. her cell farkli sector mu frekans mi bilmiyoruz, paylasmadilar. site basina cell sayisi az olan bir yer secmeye calistim. farkli combolar deneyebilirsiniz. ns3 deki LTE config setup ina bakarak. ornegin tek carrier ve her cell birer sector; veya multi carrier ve celler secotr ve carrierlara paylastirilmis vb. asagidaki ornekte sabitledim cell sayisina bakmadan 3 sectore boldum ama bakilmasi lazim}
A physical site may host multiple sectors and/or multiple LTE carriers. Build an explicit mapping \texttt{site $\rightarrow$ sector $\rightarrow$ carrier/cell}. If metadata cannot distinguish these layers, document the ambiguity and use two alternative scenarios: (A) one logical cell per observed cell ID, and (B) a capped number of physical sectors with co-sector carriers. \textbf{The classic ns-3 LTE helper does not automatically provide a faithful multi-carrier operator deployment simply because several eNB devices share a coordinate}; model carrier layers only if supported by the chosen ns-3 stack and configuration.

\begin{progress}
Hücrelerin sektör mü frekans mı olduğunu doğrudan bilmediğimiz için KPI'daki frekanslar arası handover kayıtlarına baktık. Bir hücre kullanıcılarını başka bir frekansa aktarıyorsa, sahada veya yakınında başka bir frekans var demektir. Buna göre 8 ana sahanın 5'i (130, 296, 440, 441, 443) birden fazla frekans kullanıyor; 444 ve 554 ise tek frekanslı görünüyor. 554'ün iki hücresinden biri çok yakın, diğeri çok uzak kullanıcılara hizmet veriyor; bunlar büyük olasılıkla iki farklı yöne bakan iki sektör. Hangi hücrenin hangi banda ait olduğunu ise bu veriden kesin söyleyemiyoruz.
\end{progress}


\section{First case study: S\i nd\i rg\i{} earthquake-area configuration}
\label{sec:sindirgi}
\red{istedigim regiona gore farkli dosyalardan veri cektirdim verify edin.}

\begin{progress}
Tablodaki değerleri ham Experia dosyasından tekrar hesapladık. Hücre sayılarının hepsi aynı, saha RSRP ortalamaları en fazla 0,5~dB farklı. Tek fark kayıt sayısında: seçim dosyası tam dönemden biraz daha az kayıt içeriyor, muhtemelen daha kısa bir dönemden üretilmiş. Bu sonuçları değiştirmiyor.
\end{progress}
\subsection{Geographical scope and measured inventory}
Our initial case study uses eight \emph{evaluation} sites and eight \emph{interference-buffer} sites selected from the coordinate inventory near the 10 August 2025 S\i nd\i rg\i{} earthquake epicenter (39.2461$^\circ$N, 28.0697$^\circ$E). The selection was based on proximity to the mainshock and the presence of Experia observations; it is \emph{not yet} an aftershock-footprint optimization. The epicentral coordinate is a geographical reference, not evidence of site damage. Future region selection may use an independently obtained aftershock catalog, including event coordinates, dates, and magnitudes.

\begin{table}[htbp]\centering\small
\caption{Eight selected evaluation sites. Cell counts are IDs observed in the matched Experia data, not confirmed antenna sectors. RSRP is the unweighted mean of per-cell time means and is not a coverage estimate.}
\label{tab:sindirgi_sites}
\begin{tabular}{@{}lrrrr@{}}\toprule
Site ID & Latitude & Longitude & Observed cells & Mean RSRP (dBm)\\\midrule
SITE\_000440 & 39.212651 & 28.096848 & 9 & -112.07\\
SITE\_000441 & 39.239050 & 28.164350 & 9 & -102.74\\
SITE\_000130 & 39.308051 & 27.988545 & 7 & -111.84\\
SITE\_000296 & 39.255600 & 28.183712 & 6 & -103.57\\
SITE\_000443 & 39.244750 & 28.187651 & 14 & -101.38\\
SITE\_000554 & 39.329550 & 28.149650 & 2 & -105.57\\
SITE\_000482 & 39.264157 & 28.209098 & 2 & -99.04\\
SITE\_000444 & 39.143886 & 28.158027 & 6 & -114.79\\\bottomrule
\end{tabular}
\end{table}

The evaluation set contains 55 observed cell IDs and 6,394 matched hourly radio observations. The buffer set is SITE\_000506, SITE\_000480, SITE\_000553, SITE\_000200, SITE\_000224, SITE\_000555, SITE\_000475, and SITE\_000206. These 16 physical locations are the first geographical input; the buffer is provisional and must be expanded until interior association, RSRP, and SINR statistics stabilize. The analysis files \texttt{sindirgi\_lte\_sites.csv} and \texttt{sindirgi\_lte\_cells.csv} preserve the site/cell mapping and radio summaries. A site map is supplied separately. \red{ekte paylasacagim}

\subsection{Proposed first executable topology (Configuration S0)}
Configuration S0 is a deliberately simplified, reproducible \emph{single-carrier, three-sector hypothesis} for each of the 16 physical sites: 48 logical LTE eNB sector devices in total, of which 24 are at evaluation sites. These devices are \emph{synthetic representatives}, not a one-to-one mapping to Turkcell cell IDs. The 2-cell sites are also assigned three hypothetical sectors in S0 to avoid mixing inferred and known geometry; this is an explicit modeling assumption, to be tested against one-/two-sector alternatives. No claim is made that any site actually has three sectors or uses one carrier.

\begin{enumerate}[leftmargin=*,itemsep=2pt]
\item Import all 16 coordinates and project them to local Cartesian meters; instantiate three colocated eNB devices per site, each with its own directional azimuth.
\item Assign a common illustrative LTE Band-3 carrier (DL/UL EARFCN 1300/19300), 20 MHz (100 RB), 46 dBm transmit power per modeled eNB, 30 m eNB height, 1.5 m UE height, and a 65$^\circ$ horizontal antenna beamwidth with approximately 15 dBi peak gain. These are scenario parameters, not operator configuration estimates; verify the eNB power/antenna-gain convention to avoid double-counting EIRP.
\item Set the three azimuths at site $i$ to $\phi_i$, $\phi_i+120^\circ$, and $\phi_i+240^\circ$. Draw independent $\phi_i\in[0^\circ,120^\circ)$ and repeat over an ensemble of fixed random seeds. Record the azimuth assignment with every run.
\item Use a supported LTE propagation/shadowing configuration, proportional-fair scheduling, EPC, and initially stationary UEs. Set a non-bottleneck core/backhaul baseline. Check that the chosen propagation and antenna models are supported by the pinned ns-3 release.
\item Create synthetic UEs across the evaluation area \emph{and buffer}, including uniform and spatially clustered variants. Let radio-based cell selection determine attachment where supported. Do not assign each synthetic UE to an Experia cell ID: the mapping is unknown.
\item Run a low-load radio baseline first; generate spatial maps of strongest serving-cell RSRP, SINR, serving sector, and unserved locations. Compare distributions in the interior, not just at the boundary.
\end{enumerate}

\begin{progress}
S0'ı bu adımlarla kurduk. Bilinmeyen ayarları denemeyle bulduk: güç, anten genişliği, yayılım modeli, gölgeleme ve kullanıcı yerleşimi için 49 farklı kombinasyonu ns-3'te çalıştırdık. Seçimi yalnızca 5 eğitim sahasına bakarak yaptık, kalan 3 sahayı kontrol için ayırdık. Ölçümlere en yakın ve tekrarlarda en tutarlı sonucu sektör başına 49~dBm güç, 120$^\circ$ anten genişliği ve gölgelemesiz yayılım verdi. Ayrıca 64 sanal kullanıcının yetersiz olduğunu gördük: bazı sahalara yalnızca 1--2 kullanıcı düşüyordu ve sonuçlar kullanıcıların nereye düştüğüne göre değişiyordu. Bu yüzden 256 kullanıcıya çıkardık.
\end{progress}

\paragraph{Scale and feasibility.} A first test can use 16 sites, 48 sectors, and approximately 100--200 stationary UEs; these UE counts are \emph{computational test settings}, not inferred user populations. Run a smaller 2--3-site topology first to verify antenna orientation and interference behavior, then scale to all 16 sites. A 48-sector network may be expensive with full traffic; adjust UE count and simulation duration based on measured runtime.

\subsection{Radio consistency checks and alternative sector/carrier hypotheses}
For each Turkcell cell ID, compute a time-series signature containing RSRP, SINR, RSRQ, and the \emph{raw} TAAVG value, along with sample counts and variability. Within each physical site, calculate pairwise RSRP differences, synchronized-hour correlations, and coverage of common observation hours. Correlation is descriptive evidence only: aggregate serving-user RSRP cannot identify azimuth or carrier frequency, and distinct sectors can have similar distributions.

We propose three explicit alternatives: \textbf{S0}, three synthetic single-carrier sectors at every site; \textbf{S1}, site-specific sector counts (e.g., one or two at the observed 2-cell sites, three elsewhere), selected as assumptions; and \textbf{S2}, selected sites represented by multiple carrier layers when the cell-ID inventory and radio signatures make this plausible. S2 requires correct inter-frequency modeling and should not be treated as equivalent to adding colocated cochannel sectors. Compare simulated interior RSRP/SINR distributions and service-area geometry under S0/S1/S2. Report whether conclusions persist across hypotheses; do not label an inferred configuration as confirmed.

\begin{progress}
Üç senaryoyu da denedik. İki hücreli sahalarda sektör sayısını azaltmak (S1) bir iyileşme sağlamadı. Çok frekanslı senaryoda (S2), handover kayıtlarına göre çok frekanslı görünen sahalara 1800~MHz'e ek olarak bir 800~MHz katmanı ekledik. S2 hücrelerin sinyal düzeyini ve CQI'yı gerçeğe yaklaştırdı, ama saha ortalamalarındaki hatayı artırdı (kontrol sahalarında RSRP farkı 1,7'den 3,0~dB'ye çıktı). Bunun nedeni, 800~MHz katmanının gücünü ve kullanıcıların frekanslara nasıl dağıldığını bilmememiz. Bu yüzden şimdilik tek frekanslı S0'ı temel senaryo olarak kullanıyoruz.

Yük deneylerinde S0 ve S2'yi birlikte koşturmayı öneriyoruz. S0'da girişim olduğundan fazla, S2'de az; iki senaryoda da aynı çıkan sonuç daha güvenilir olur. Bant bilgisi gelirse S2 netleşir.
\end{progress}

\paragraph{Experia timing-advance caveat.} \red{bunu Turkcell le netlestirmeye calisiyoruz. bazi degerler anlamsiz}The supplied vendor description defines 12 cell-range usage categories: 0--9\%, 10--19\%, \ldots, 100--109\%, and 110--200\%. However, the exported \texttt{AVG(TAAVG)} column contains values far outside the 0--11 category-index range. Therefore, neither $\mathrm{TAAVG}\times 1.5$ km nor a standard LTE timing-advance-command conversion is justified for the raw export. Preserve raw values and exclude TAAVG from distance-based calibration until its aggregation, scaling, and units are verified. Even with a valid conversion, the metric describes weighted usage distance, not maximum cell radius or UE coordinates.

\begin{progress}
TAAVG'yi birimi belli olmadığı için mesafeye çevirmedik. Yine de göreli olarak anlamlı görünüyor: aynı sahada TAAVG'si büyük olan hücrelerin sinyali genelde daha zayıf ve KPI'daki RACH TA da aynı yönü gösteriyor. Birim netleşirse mesafe kontrolüne eklenebilir.
\end{progress}

\subsection{Staged calibration and acceptance tests}
\begin{enumerate}[leftmargin=*,itemsep=2pt]
\item \textbf{Data audit:} verify cell-ID joins, missing hourly measurements, time coverage, and outliers; record per-site and per-cell counts. Pay special attention to SITE\_000444, which has strongly uneven observation counts across its cells.
\item \textbf{Geometry sanity:} verify projected inter-site distances, directional antenna patterns, cochannel interference, and association; plot RSRP and SINR grids.
\item \textbf{Distributional radio check:} compare measured cell-level RSRP/SINR/RSRQ summaries with simulated \emph{serving-UE} summaries under explicit UE sampling. Because the observed UE distribution is unknown, evaluate multiple synthetic placements and report uncertainty; do not fit only regional means.
\item \textbf{Holdout:} reserve selected cells or time windows for checks not used to tune propagation/shadowing and UE placement. Since simulation sectors do not map directly to operator cells, prioritize regional distribution and site-level range checks over unjustified one-to-one cell matching.
\item \textbf{Buffer convergence:} add more neighboring sites until interior SINR, RSRP, and serving-cell assignments change negligibly under a predeclared tolerance.
\item \textbf{Traffic stage:} only after radio plausibility checks, introduce synthetic applications and calibrate broad throughput/PRB distributions against the separate, non-contemporaneous traffic dataset.
\end{enumerate}

\begin{progress}
Bu adımların hepsini yaptık. Ayarları 5 eğitim sahasıyla seçip 3 kontrol sahasında sınadık; komşu sayısı testi de geçti. Simülasyon, hangi sahanın güçlü hangisinin zayıf olduğunu doğru sıralıyor. Ancak hücre düzeyinde bakınca, gerçek ağdaki çok güçlü hücreler simülasyonda yok ve simülasyondaki hücreler birbirine fazla benziyor. Bunun nedeni büyük olasılıkla gerçek ağın çok frekanslı olması ve gerçek kullanıcıların sahalara bizim varsaydığımızdan daha yakın olması.
\end{progress}

\subsection{Portable region-generation interface}
The same code should accept \texttt{region.yaml} (geographic polygon or epicenter/aftershock selection, evaluation sites, guard-ring rule, assumed carrier/sector policy, UE spatial model, random seeds) plus \texttt{sites.csv} and \texttt{cell\_radio.csv}. The S\i nd\i rg\i{} case is the first configuration; coastal and city-center scenarios should reuse the exact generator with different geographic subsets and environment assumptions. This separates geographical input, uncertain radio configuration, and experimental demand/failure scenarios.

\begin{progress}
Henüz yapılmadı. Şu an bölge ayarları her deneyin kendi dosyalarında duruyor. Şehir merkezi veya sahil bölgesine geçmeden önce bunları tek bir bölge dosyasında toplamak gerekecek.
\end{progress}

\section{Baseline radio configuration}
\red{Turkcell radio ve frekans verisi paylasmadi. o yuzden nes3 default degerleri ile basliyoruz. handover vs de tanimli ama en bast astatic UE yapabiliriz}
Table~\ref{tab:radio} provides \emph{assumed} starting values, not inferred Turkcell settings. Check support and attribute names against the installed ns-3 release before implementation.
\begin{table}[ht]\centering\small
\caption{Representative LTE macro-cell baseline.}\label{tab:radio}
\begin{tabularx}{\linewidth}{@{}l l X@{}}\toprule
Parameter & Initial value & Interpretation / caveat \\
\midrule
Carrier & LTE Band 3, 1800 MHz & Scenario assumption; choose consistent DL/UL EARFCNs \\
DL / UL EARFCN & 1300 / 19300 & Illustrative paired Band-3 values (DL 1815 MHz, UL 1720 MHz) \\
Bandwidth & 20 MHz (100 RB) & LTE PHY bandwidth in resource blocks \\
eNB power & 46 dBm & Per modeled eNB carrier/sector; check reference-signal power mapping \\
eNB / UE height & 30 m / 1.5 m & Representative macro / pedestrian \\
UE maximum TX power & 23 dBm & Representative UE limit \\
Sector count & up to 3/site & Only if supported by cell metadata \\
Horizontal pattern & 65$^\circ$ beamwidth & Directional model with approx. 15 dBi maximum gain \\
Azimuth & $\phi_i+120^\circ s$ & Independent unknown rotation $\phi_i$ per 3-sector site \\
Downtilt & 6$^\circ$ if supported & Do not claim 3D pattern if using horizontal-only antenna \\
Propagation & urban macro candidate & Compare supported models; record all parameters \\
Shadowing & 8 dB std. dev. & Assumed; assess sensitivity and spatial correlation \\
Scheduler & proportional fair & \texttt{PfFfMacScheduler} \\
Handover & A3 RSRP-based & Appropriate for mobile/failure experiments; not load-aware \\
EPC / backhaul & non-bottleneck baseline & Keep transport effects separate from radio effects \\
\bottomrule\end{tabularx}\end{table}

\begin{progress}
ns-3.44'ün klasik LTE modülünü kullandık. Tablodaki değerlerle başladık, sonra ölçümlere göre değiştirdik. Seçtiğimiz değerler taslaktakilerden biraz farklı:
\begin{center}\small
\begin{tabular}{@{}lll@{}}\toprule
 & Taslak & Kullandığımız \\\midrule
Güç & 46~dBm & 49~dBm \\
Anten genişliği & 65$^\circ$ & 120$^\circ$ \\
Yayılım & urban macro & log-distance (üs 3) \\
Gölgeleme & 8~dB & yok \\
Sanal kullanıcı & 100--200 & 256 \\
Yük & --- & KPI'daki ortalama kullanım \\\bottomrule
\end{tabular}
\end{center}
Bunlar Turkcell'in gerçek ayarları değil; denediklerimiz arasında ölçümlere en yakın ve tekrarlarda en tutarlı sonucu verenler. Kullanıcılar sabit olduğu için handover kullanmadık.
\end{progress}

\subsection{Azimuth uncertainty}
For site $i$ with three sectors, define $\theta_{i,s}=\phi_i+120^\circ s$, $s=0,1,2$, and sample $\phi_i\sim\mathrm{Uniform}(0^\circ,120^\circ)$ independently across sites. For two sectors use a separate assumption (e.g., 180$^\circ$ separation). Keep random seeds and azimuth realizations fixed when comparing algorithms. Sector patterns and association boundaries are outputs of the radio simulation, not predetermined circular coverage polygons.

\section{Propagation, RSRP, SINR, and cell-size plausibility}
\label{sec:prop}
For candidate serving cell $i$ and location $x$, the reference-signal received power depends on reference-signal power, antenna pattern, path loss, and shadowing:
\begin{equation}
P_{\mathrm{RSRP},i}(x)\approx P_{\mathrm{RS},i}+G_i(x)-PL_i(x)-L_i(x).
\end{equation}
SINR further depends on neighboring transmissions and interference, so matching RSRP alone is insufficient. Select one supported urban-macro propagation model for the baseline; compare at least one alternate model. Explicitly report whether shadowing, fading, building losses, and interference are enabled. The chosen LTE stack and ns-3 version determine which standardized models and 3D antenna features are available.

\paragraph{Use of timing advance.} Mean timing advance may constrain approximate serving-UE distance. LTE's basic timing-advance command step corresponds to about 78 m one-way range under the standard conversion; a mean alone does not identify the distance distribution, azimuth, or cell boundary. Don't infer actual UE positions from TAAVG alone.

\paragraph{Radio plausibility checks.} Generate a grid of candidate UE positions, calculate strongest-cell association and RSRP/SINR maps, and compare distributions of simulated serving-cell RSRP, SINR, and TA-derived distances with radio observations. Since observations and traffic are from different periods, treat these as distributional plausibility checks, not time-aligned validation. RSRQ additionally depends on load and interference and should not be forced to match a traffic period that differs from its measurement period.

\begin{progress}
Bölgenin 100 metre çözünürlükte kapsama haritalarını çıkardık: hangi sahanın nereye hizmet verdiği, sinyal gücü ve sinyal kalitesi. Haritaların ns-3 ile birebir aynı sonucu verdiğini kontrol ettik. Değerlendirme alanının \%99,8'inde sinyal yeterli; yani bölgede kapsama boşluğu yok. Tek frekans varsayımında ağın asıl sorunu sinyal zayıflığı değil, komşu sahaların girişimi: bütün komşular aynı anda yayın yaptığında alanın yarısından fazlasında sinyal kalitesi düşüyor.

Bir uyarı: ns-3'ün raporladığı SINR değerlerinin çoğu, ns-3'ün hesaplama biçimi nedeniyle girişimi içermiyor. Ayrıca ns-3 RSRQ'yu yükten bağımsız hesaplıyor. Bu yüzden SINR karşılaştırmalarını temkinli yorumluyoruz, RSRQ'yu ise karşılaştırmada kullanmıyoruz.
\end{progress}

\section{Synthetic UEs and initial association}
Generate a configurable number of UEs in the selected region. Compare (U1) uniform spatial placement, (U2) clustered hotspots, and (U3) radial placement constrained by cell-level TA statistics when meaningful. A spatial UE density is a scenario assumption, not observed ground truth. Start with stationary UEs to isolate coverage and interference; add mobility later.

In the classic LTE helper, \texttt{Attach(ueDevices)} can perform automatic initial cell selection based on radio measurements, whereas \texttt{Attach(ueDevice,enbDevice)} forces a chosen eNB. Verify the exact overload and handover behavior in the installed version. Configure X2 interfaces and an A3 RSRP handover algorithm for mobility studies; neither RSRP-based handover nor initial attachment automatically implements load balancing. For site outages or UAV eNB introduction, explicit recovery/reselection logic may be needed.

\begin{progress}
Kullanıcıları hem bölgeye düzgün dağıtarak hem de sahaların çevresinde kümeler hâlinde denedik. En tutarlı sonuç sahaların 2--6~km çevresindeki kümelerle çıktı. Kullanıcılar şimdilik sabit ve handover henüz kurulmadı; kesinti ve UAV deneylerinde gerekecek.
\end{progress}

\section{Traffic model and calibration}
\label{sec:traffic}
For KPI interval duration $\Delta t$, convert carried traffic volume $V_{i,t}$ into measured mean throughput $\bar S_{i,t}=8V_{i,t}/\Delta t$ (with unit conversion). This is \emph{served traffic}, not necessarily offered demand. Begin with synthetic downlink/uplink UDP sources, calibrated so uncongested simulations produce similar per-cell carried traffic and broad PRB utilization patterns. The exact application types, packet sizes, arrival bursts, and per-UE allocations remain experimental choices. Use active-user counts as constraints on synthetic UE populations if their KPI definitions are suitable.

\paragraph{Avoid circular validation.} Reserve some observables (e.g., CQI or per-user throughput distributions) and some cells/times for out-of-sample checks rather than tuning all parameters to every KPI. Because actual bandwidth and scheduler are unknown, matching PRB utilization may require a scenario-specific demand scaling and cannot uniquely recover capacity. Reproduce both distributions and spatial variation; report errors by cell and across cells.

\begin{progress}
Bu bölümden üç önemli bulgu çıktı.

\textbf{Simülasyondaki ağ çok boştu.} KPI'ya göre gerçek ağda kapasitenin normal günde \%23'ü, akşam yoğun saatinde \%36'sı kullanılıyor. İlk simülasyonumuzda bu oran yalnızca \%3'tü; yani ağı gerçeğinden 7--10 kat boş çalıştırıyorduk. Her hücreye trafik üreten bir kullanıcı ekleyerek gerçekçi yükü verdik. Bunun sonucunda, daha önce olduğundan iyi görünen sinyal kalitesi (SINR) ölçümlere belirgin şekilde yaklaştı.

\textbf{Deprem anı KPI'da açıkça görülüyor.} 10 Ağustos ve 27 Ekim depremleri, KPI verisinde depremin olduğu 15 dakikalık dilimde fark ediliyor. Depremden sonraki ilk yarım saatte aynı anda bağlı kullanıcı sayısı yaklaşık iki katına çıkıyor, indirilen veri miktarı ise yaklaşık üçte bir azalıyor ve bağlantı kopmaları artıyor. Yani deprem yükü ``daha çok veri'' değil, ``aynı anda çok daha fazla kişi''. Stres testini yalnızca trafik miktarını artırarak değil, aynı anda bağlı kişi sayısını artırarak da kurmayı öneriyoruz.

\textbf{Bağımsız kontrol olarak CQI.} Yükü ayarlarken CQI'yı hiç kullanmadık, sonra bağımsız bir kontrol olarak karşılaştırdık. KPI'da CQI ve yük aynı kayıttan geldiği için bu, aynı dönemin verisiyle yapılan bir karşılaştırma. Simülasyonun CQI'sı gerçek ağa yakın ama ortalamada yaklaşık 1 puan düşük. Bu da simülasyonda girişimin biraz fazla olduğunu, yani tek frekans varsayımının girişimi abarttığını gösteriyor.
\end{progress}

\section{Implementation sequence in ns-3}
\begin{enumerate}[leftmargin=*,itemsep=3pt]
\item \textbf{Pin a simulator version.} Record ns-3 release, LTE/EPC module version, compiler, build flags, and available propagation/antenna models. Use a small stock LTE+EPC example as a smoke test.
\item \textbf{Import topology.} Read \texttt{sites.csv}, convert lat/lon to local meters, create eNB nodes, and install constant-position mobility with a representative antenna height.
\item \textbf{Build sector/cell devices.} Instantiate logical eNB devices per chosen sector/carrier mapping; set unique cell IDs internally, radio attributes, and directional antenna orientations. Verify that colocated sectors are independent devices where intended.
\item \textbf{Configure LTE and EPC.} Set valid paired EARFCNs, RB bandwidth, TX power, propagation and fading models, scheduler, EPC, and remote-host routes. Ensure interference is represented across cells operating on the same carrier.
\item \textbf{Create UEs.} Generate UE locations with reproducible seeds, install UE LTE devices, configure initial cell selection, and verify that every UE has a sensible serving cell and IP route.
\item \textbf{Produce unloaded radio maps.} Export grid-based RSRP/SINR and serving-cell assignments. Plot association areas and inspect boundary artifacts. Do not equate strongest-RSRP areas with QoS service areas.
\item \textbf{Add baseline traffic.} Configure DL/UL applications and collect per-cell carried traffic, UE throughput, PRB usage, CQI, RLC/PDCP counters, and handover events as supported.
\item \textbf{Compare with observations.} Check per-cell and network-level distributions of RSRP/TA (radio period) and traffic/PRB/CQI/throughput (traffic period), without combining mismatched timestamps.
\item \textbf{Generate an ensemble.} Repeat over independent sector rotations, UE distributions, shadowing seeds, and selected radio configurations. Retain plausible scenarios; report which conclusions are sensitive to assumptions.
\item \textbf{Freeze a baseline.} Export complete configuration files, random seeds, per-cell mappings, UE locations, maps, and validation plots for later stress tests.
\end{enumerate}

\subsection{Illustrative ns-3 configuration fragment}
The following is an \emph{illustrative fragment}, not a standalone program; attribute support should be verified for the pinned release. In particular, antenna orientation is set per eNB device/site, not globally if each site has a different rotation.
\begin{lstlisting}[style=cpp]
Ptr<LteHelper> lte = CreateObject<LteHelper>();
Ptr<PointToPointEpcHelper> epc =
    CreateObject<PointToPointEpcHelper>();
lte->SetEpcHelper(epc);
lte->SetEnbDeviceAttribute("DlEarfcn", UintegerValue(1300));
lte->SetEnbDeviceAttribute("UlEarfcn", UintegerValue(19300));
lte->SetEnbDeviceAttribute("DlBandwidth", UintegerValue(100));
lte->SetEnbDeviceAttribute("UlBandwidth", UintegerValue(100));
lte->SetEnbPhyAttribute("TxPower", DoubleValue(46.0));
lte->SetSchedulerType("ns3::PfFfMacScheduler");
lte->SetHandoverAlgorithmType(
    "ns3::A3RsrpHandoverAlgorithm");
// Next: select supported path-loss and antenna models;
// create positioned eNB and UE nodes; install devices;
// configure per-sector azimuths; add X2, IP and traffic.
\end{lstlisting}

\section{Verification, validation, and acceptance criteria}
\begin{table}[ht]\centering\small
\begin{tabularx}{\linewidth}{@{}l X X@{}}\toprule
Level & Test & Interpretation \\
\midrule
Software verification & Correct eNB coordinates, device counts, EARFCN/bandwidth, attachment, routing, traffic accounting & Required before any research conclusions \\
Radio plausibility & RSRP/SINR and TA summary distributions within defensible tolerances & Not proof of actual cell boundaries \\
Traffic plausibility & Cell-level carried traffic, utilization, CQI/throughput distributions & Depends on unknown scheduler/bandwidth and UE placement \\
Robustness & Repeat under antenna, UE, propagation, and bandwidth uncertainty & Reveals scenario dependence \\
Out-of-sample & Hold out sites and/or intervals where available & Tests limited generalization \\
\bottomrule\end{tabularx}\end{table}
Do not specify numerical acceptance thresholds before auditing KPI units, spatial variability, and model capability. Use distributional error (e.g., quantile differences), site-level errors, and stability of spatial rankings; avoid relying on one aggregate mean. A simulated network can be useful for controlled counterfactual experiments without being validated as a faithful reconstruction of Turkcell's real deployment.

\begin{progress}
Sonuçları tek bir ortalamayla değil; saha hataları, hücre dağılımları ve sahaların sıralamasıyla birlikte değerlendirdik. Her deneyin başarı ölçütünü koşudan önce yazdık; 64 kullanıcının yetersiz olduğu da böyle bir ölçütle ortaya çıktı. Kabul edilebilir hata düzeyi konusunda görüşünüzü almak istiyoruz.
\end{progress}

\section{Deliverables and staged milestones}
\begin{enumerate}[nosep]
\item \textbf{M1: Input audit.} Cell/site/carrier inventory, coordinate map, KPI dictionary, and mismatch report.
\item \textbf{M2: Geometry-only network.} Reproducible eNB/sector placement and radio environment maps.
\item \textbf{M3: UE association and baseline.} UE distribution scenarios, attachment checks, LTE/EPC traffic, and baseline KPIs.
\item \textbf{M4: Measurement plausibility.} RSRP/TA and traffic/KPI comparisons, plus uncertainty ensemble.
\item \textbf{M5: Frozen scenario generator.} Configuration/seed manifest and reusable ns-3 scripts for subsequent studies.
\end{enumerate}

\begin{progress}
M1--M4 tamamlandı: veri denetimi, kapsama haritaları, kullanıcı ve yük senaryoları ve ölçümlerle karşılaştırma. M5 kısmen hazır; senaryo ve kodlar tekrar kullanılabilir, ama başka bölgelere taşımak için tek bir bölge dosyası henüz yok.
\end{progress}

\section{Future extensions }
\paragraph{Stress testing.} Increase offered load globally or in spatial hotspots; separately introduce site outages or backhaul limits. Define \emph{service failure} using throughput, delay, drop, or unmet-demand thresholds, not physical equipment failure inferred from congestion. Evaluate UE reassociation and spillover congestion.

\begin{progress}
Stres testine henüz başlamadık. Önerimiz, yoğun saat yükünden başlayıp iki şeyi ayrı ayrı adım adım artırmak: ağdaki trafik miktarı ve deprem anını temsil etmek için aynı anda bağlı kullanıcı sayısı. Her iki deneyi tek ve çok frekanslı senaryoda koşturup hizmetin nerede yetersiz kaldığını (hız, gecikme, bağlanamayan kullanıcı) karşılaştırabiliriz.
\end{progress}
\paragraph{UAV-assisted restoration.} Translate simulated unmet demand into geographical service targets with dynamic priorities and deadlines. Add UAV eNBs with explicit access capacity, backhaul and interference; evaluate actual UE reassociation and service restoration. Train priority-window RL in a faster surrogate environment and validate resulting deployments in ns-3. This is a separate later-stage research question.

\section{Open decisions before coding}
\begin{enumerate}[nosep]
\item Which ns-3 release and LTE stack are available? Are 3D antenna patterns and the desired propagation model supported?
\item Does each observed cell ID denote a sector, a carrier, or both? Can any metadata distinguish bands?
\item What are the exact definitions/units of TAAVG, traffic volume, throughput, PRB utilization, and active users?
\item How many days and cells are covered by each dataset, and what fraction have matching cell IDs?
\item Which geographical cluster offers enough surrounding sites to limit edge effects?
\end{enumerate}

\begin{progress}
Sizden ve Turkcell'den öğrenmek istediklerimiz:
\begin{itemize}[nosep]
\item Hücrelerin frekans bandı ve gücü, en azından hangi hücrelerin aynı sektörde olduğu. Bu bilgi gelirse çok frekanslı senaryoyu tahmine değil veriye dayandırabiliriz.
\item Experia'daki SINR ve RSRQ'nun hangi kanal ve ortalamayla hesaplandığı. Simülasyonla aynı şeyi karşılaştırdığımızdan emin olmak istiyoruz.
\item KPI'daki throughput ve TAAVG birimleri. Bunlar netleşirse yükü ve mesafeyi daha doğru modelleyebiliriz.
\item KPI saatlerinin yerel saat olduğunun teyidi. Deprem anındaki sıçramalar buna işaret ediyor.
\item Bu başlangıç senaryosunun, sınırlarıyla birlikte, stres testine geçmek için yeterli olup olmadığı.
\end{itemize}
\end{progress}

\paragraph{Implementation reference.} Consult the ns-3 LTE user documentation and the examples distributed with the \emph{installed version}: \url{https://www.nsnam.org/docs/models/html/lte-user.html}. API and model support may differ between releases.
\end{document}
